\begin{problem}[2020第37届全国中学生物理竞赛复赛][回旋加速器原理与反质子产生]{120}
    劳伦斯（E. O. Lawrence）在1930年首次提出了回旋加速器的原理：用两个半圆形磁场，使带电粒子沿圆弧形轨道旋转，反复通过两半圆缝隙间的高频电场加速而获得较高能量。他因这个极富创意的方案而获得了1939年的诺贝尔物理学奖。
    \begin{subproblem}
        目前全球最大的回旋加速器是费米实验室中的高能质子同步加速器 Tevatron（粒子运行最大回旋圆轨道的周长为 $L_{\max} = 6436 \mathrm{~m}$），可以将一质子加速到的最大能量为 $E_{\mathrm{p,max}}^{(\mathrm{tevatron})} = 1.00 \times 10^{6} \mathrm{MeV}$。假设质子在被加速过程中始终在垂直于均匀磁场的平面内运动，不计电磁辐射引起的能量损失。求该同步加速器 Tevatron 将质子加速到上述最大能量所需要的磁感应强度的最小值 $B_{\min}$。
    \end{subproblem}
    \begin{subproblem}
        高能入射质子轰击静止的质子（靶质子），可产生反质子 $\overline{\mathrm{p}}$，反应式为
        \begin{equation}
            \mathrm{p} + \mathrm{p} \rightarrow \mathrm{p} + \mathrm{p} + \mathrm{p} + \overline{\mathrm{p}}
        \end{equation}
        求能产生反质子时入射质子的最小动能，并判断第 (1) 问中的 Tevatron 加速的质子是否可以轰击静止的靶质子而产生反质子。

        已知数据：质子质量 $m_{\mathrm{p}} = 938.3 \mathrm{MeV} / c^{2}$，真空中的光速 $c = 3.00 \times 10^{8} \mathrm{m/s}$。
    \end{subproblem}
\end{problem}